%---------------------------------------------------------------------
% UNIT 2 --- RANDOM VARIABLES
% Source: Section 2.7 Self Assessment Questions, pp. 42-43.
%---------------------------------------------------------------------
\chapter{Random Variables}

\srcnote{Source: \emph{Business Mathematics} 5405/1429, \S2.7 \saq{Self Assessment Questions}, pp.~42--43.}

\begin{keyidea}[What makes a list of numbers a probability distribution]
A function $P(X=x)$ is a discrete probability distribution if and only if
\[
0\le P(X=x)\le 1 \quad\text{for every }x,
\qquad\text{and}\qquad
\sum_{\text{all }x}P(X=x)=1 .
\]
The second condition is the whole of Questions 8, 9 and 10, and it is the check
that catches an arithmetic slip in Questions 3 to 7.
\end{keyidea}

%---------------------------------------------------------------------
\section{Defining a random variable}

\begin{question}{Q.1}
Define random variable. What is the difference between a discrete random
variable and a continuous random variable?
\end{question}

\begin{solution}
A \textbf{random variable} is a rule that assigns one real number to each
outcome of a random experiment. Formally it is a function $X:S\to\mathbb{R}$,
where $S$ is the sample space. It converts outcomes that may not be numeric
(``head'', ``defective'') into numbers that can be averaged, added and
distributed.

\begin{center}
\small\sffamily
\begin{tabular}{@{}L{3.0cm} L{5.6cm} L{5.6cm}@{}}
\toprule
 & \textbf{Discrete} & \textbf{Continuous} \\
\midrule
Values taken & A countable set --- typically integers & Every value in an interval \\
Arises from & Counting & Measuring \\
Probability of a point & $P(X=x)$ can be positive & $P(X=x)=0$ always \\
Described by & Probability mass function & Probability density function $f(x)$ \\
Total probability & $\sum_x P(X=x)=1$ & $\int_{-\infty}^{\infty} f(x)\,\mathrm{d}x=1$ \\
Business example & Number of defective units in a batch & Time to serve a customer \\
\bottomrule
\end{tabular}
\end{center}

\ans{A random variable assigns a real number to each outcome of an experiment.
It is discrete if its values are countable (results of counting) and continuous
if it can take any value in an interval (results of measuring). For a continuous
variable $P(X=x)=0$ for every single point, so probability is found over
intervals via a density function.}
\end{solution}

%---------------------------------------------------------------------
\section{Discrete random variables in business}

\begin{question}{Q.2}
Write down the applications of the discrete random variable in business
mathematics.
\end{question}

\begin{solution}
Any business quantity that is obtained by \emph{counting} is a discrete random
variable, and each supports a standard decision:

\begin{description}
  \item[Quality control] The number of defective items in a sample of $n$ is
        binomial. It drives accept/reject rules for incoming batches.
  \item[Demand and inventory] Units of a product sold per day is discrete, and
        its distribution sets the re-order level and the safety stock.
  \item[Queue arrivals] Customers arriving at a counter per hour is typically
        Poisson, which fixes how many tellers to staff.
  \item[Insurance and warranty] The number of claims in a period is discrete;
        its expected value is the premium loading.
  \item[Credit and default] The number of accounts defaulting out of $n$ loans is
        binomial, and drives the provision held against the portfolio.
  \item[Expected monetary value] For any of the above,
        $E(X)=\sum x\,P(X=x)$ turns the distribution into the single number a
        manager decides on.
\end{description}

\ans{Quality control (defectives per batch), demand and inventory planning,
arrivals in queueing and staffing, insurance claim counts, and default counts in
credit risk --- each summarised for decision-making by
$E(X)=\sum x\,P(X=x)$.}
\end{solution}

%---------------------------------------------------------------------
\section{False alarms: frequency to probability}

\begin{question}{Q.3}
The fire chief for a small volunteer fire department has compiled data on the
number of false alarms called in each day for the past 360 days. The frequency
distribution below summarises the findings. Construct the probability
distribution for this study.

\medskip
\centering\small\sffamily
\begin{tabular}{@{}C{3.2cm} C{2.4cm}@{}}
\toprule
\textbf{No.\ of false alarms} & \textbf{Frequency} \\
\midrule
0 & 75 \\ 1 & 80 \\ 2 & 77 \\ 3 & 40 \\
4 & 28 \\ 5 & 24 \\ 6 & 20 \\ 7 & 16 \\
\midrule
\textbf{Total} & \textbf{360} \\
\bottomrule
\end{tabular}
\end{question}

\begin{solution}
Converting a frequency distribution into a probability distribution is one
division, applied to every row: relative frequency estimates probability,
\[
P(X=x)=\frac{f(x)}{\sum f}=\frac{f(x)}{360}.
\]

\begin{center}
\small\sffamily
\begin{tabular}{@{}C{2.0cm} C{2.0cm} C{2.6cm} C{2.4cm}@{}}
\toprule
$x$ & $f(x)$ & $P(X=x)=f/360$ & Decimal \\
\midrule
0 & 75  & $75/360$ & 0.2083 \\
1 & 80  & $80/360$ & 0.2222 \\
2 & 77  & $77/360$ & 0.2139 \\
3 & 40  & $40/360$ & 0.1111 \\
4 & 28  & $28/360$ & 0.0778 \\
5 & 24  & $24/360$ & 0.0667 \\
6 & 20  & $20/360$ & 0.0556 \\
7 & 16  & $16/360$ & 0.0444 \\
\midrule
\textbf{Total} & \textbf{360} & \textbf{1} & \textbf{1.0000} \\
\bottomrule
\end{tabular}
\end{center}

\vcheck{$75+80+77+40+28+24+20+16=360$, so the probabilities sum to
$360/360=1$ exactly. The distribution is valid.}

\medskip
As a by-product, the expected number of false alarms per day is
\[
E(X)=\sum x\,P(X=x)
=\frac{0(75)+1(80)+2(77)+3(40)+4(28)+5(24)+6(20)+7(16)}{360}
\]
\[
=\frac{818}{360}\approx 2.27 \ \text{false alarms per day}.
\]

\ans{$P(X=x)=f(x)/360$ for $x=0,\dots,7$, giving $0.2083$, $0.2222$, $0.2139$,
$0.1111$, $0.0778$, $0.0667$, $0.0556$, $0.0444$, which sum to 1. Mean
$\approx 2.27$ false alarms per day.}
\end{solution}

%---------------------------------------------------------------------
\section{Three tosses of a fair coin}

\begin{question}{Q.4}
Construct the discrete probability distribution which corresponds to the
experiment of tossing a fair coin three times. Let $X$ equal the number of heads
occurring in three tosses. What is the probability of two or more heads?
\end{question}

\begin{solution}
There are $2^{3}=8$ equally likely outcomes. Listing them makes the counts
undeniable:

\begin{center}
\small\sffamily
\begin{tabular}{@{}C{2.0cm} L{5.6cm} C{2.0cm} C{2.0cm}@{}}
\toprule
$x$ (heads) & Outcomes & Count & $P(X=x)$ \\
\midrule
0 & TTT                & 1 & $1/8$ \\
1 & HTT, THT, TTH      & 3 & $3/8$ \\
2 & HHT, HTH, THH      & 3 & $3/8$ \\
3 & HHH                & 1 & $1/8$ \\
\midrule
\textbf{Total} & & \textbf{8} & \textbf{1} \\
\bottomrule
\end{tabular}
\end{center}

Equivalently, $X\sim\text{Binomial}(n=3,\;p=\tfrac12)$ and
$P(X=x)=\binom{3}{x}\left(\tfrac12\right)^{3}$.

\medskip
Two or more heads means $X=2$ or $X=3$, which are mutually exclusive:
\[
P(X\ge 2)=P(2)+P(3)=\frac{3}{8}+\frac{1}{8}=\frac{4}{8}=\frac{1}{2}=0.5 .
\]

\ans{$P(0)=1/8$, $P(1)=3/8$, $P(2)=3/8$, $P(3)=1/8$; and
$P(X\ge 2)=1/2=0.5$.}
\end{solution}

%---------------------------------------------------------------------
\section{Four balls drawn without replacement}

\begin{question}{Q.5}
A bag contains 4 red and 6 black balls. A sample of 4 balls is selected from the
bag without replacement. Let $X$ be the number of red balls. Find the
probability distribution of $X$.
\end{question}

\begin{solution}
Sampling \emph{without} replacement from a finite bag makes this hypergeometric,
not binomial. With 4 red and 6 black in a bag of 10, drawing 4:
\[
P(X=x)=\frac{\dbinom{4}{x}\dbinom{6}{4-x}}{\dbinom{10}{4}},\qquad x=0,1,2,3,4 .
\]
The denominator is $\binom{10}{4}=210$ for every row.

\begin{center}
\small\sffamily
\begin{tabular}{@{}C{1.4cm} L{4.6cm} C{2.0cm} C{2.2cm} C{1.8cm}@{}}
\toprule
$x$ & Numerator & Value & $P(X=x)$ & Decimal \\
\midrule
0 & $\binom{4}{0}\binom{6}{4}=1\times 15$ & 15 & $15/210=1/14$   & 0.0714 \\
1 & $\binom{4}{1}\binom{6}{3}=4\times 20$ & 80 & $80/210=8/21$   & 0.3810 \\
2 & $\binom{4}{2}\binom{6}{2}=6\times 15$ & 90 & $90/210=3/7$    & 0.4286 \\
3 & $\binom{4}{3}\binom{6}{1}=4\times 6$  & 24 & $24/210=4/35$   & 0.1143 \\
4 & $\binom{4}{4}\binom{6}{0}=1\times 1$  & 1  & $1/210$         & 0.0048 \\
\midrule
\textbf{Total} & & \textbf{210} & \textbf{1} & \textbf{1.0000} \\
\bottomrule
\end{tabular}
\end{center}

\vcheck{$15+80+90+24+1=210=\binom{10}{4}$, so the probabilities sum to 1.}

\ans{$P(0)=\tfrac{1}{14}$, $P(1)=\tfrac{8}{21}$, $P(2)=\tfrac{3}{7}$,
$P(3)=\tfrac{4}{35}$, $P(4)=\tfrac{1}{210}$.}
\end{solution}

\begin{pitfall}
Treating this as binomial with $p=4/10$ gives $P(2)=0.3456$ instead of
$0.4286$. Binomial requires a constant $p$, which only holds if the ball is
replaced. ``Without replacement'' means hypergeometric.
\end{pitfall}

%---------------------------------------------------------------------
\section{Three balls from five white and three black}

\begin{question}{Q.6}
Three balls are drawn from a bag containing 5 white and 3 black balls. If $X$
represents the number of white balls drawn, find the probability distribution of
$X$.
\end{question}

\begin{solution}
Same structure as Q.5, with $\binom{8}{3}=56$ as the denominator:
\[
P(X=x)=\frac{\dbinom{5}{x}\dbinom{3}{3-x}}{\dbinom{8}{3}},\qquad x=0,1,2,3 .
\]

\begin{center}
\small\sffamily
\begin{tabular}{@{}C{1.4cm} L{4.6cm} C{2.0cm} C{2.2cm} C{1.8cm}@{}}
\toprule
$x$ & Numerator & Value & $P(X=x)$ & Decimal \\
\midrule
0 & $\binom{5}{0}\binom{3}{3}=1\times 1$  & 1  & $1/56$          & 0.0179 \\
1 & $\binom{5}{1}\binom{3}{2}=5\times 3$  & 15 & $15/56$         & 0.2679 \\
2 & $\binom{5}{2}\binom{3}{1}=10\times 3$ & 30 & $30/56=15/28$   & 0.5357 \\
3 & $\binom{5}{3}\binom{3}{0}=10\times 1$ & 10 & $10/56=5/28$    & 0.1786 \\
\midrule
\textbf{Total} & & \textbf{56} & \textbf{1} & \textbf{1.0000} \\
\bottomrule
\end{tabular}
\end{center}

\vcheck{$1+15+30+10=56=\binom{8}{3}$.}

\ans{$P(0)=\tfrac{1}{56}$, $P(1)=\tfrac{15}{56}$, $P(2)=\tfrac{15}{28}$,
$P(3)=\tfrac{5}{28}$.}
\end{solution}

%---------------------------------------------------------------------
\section{Five rolls of a die}

\begin{question}{Q.7}
A fair die is rolled 5 times. Let $X$ represent the number of times the face 3
turns up. Obtain the probability distribution of $X$.
\end{question}

\begin{solution}
Here the rolls \emph{are} independent and $p=\tfrac16$ does not change, so this
is binomial --- the contrast with Questions 5 and 6 is deliberate:
\[
X\sim\text{Binomial}\!\left(n=5,\;p=\tfrac16\right),\qquad
P(X=x)=\binom{5}{x}\left(\frac{1}{6}\right)^{x}\left(\frac{5}{6}\right)^{5-x}.
\]
Every term has denominator $6^{5}=7776$.

\begin{center}
\small\sffamily
\begin{tabular}{@{}C{1.4cm} L{5.4cm} C{2.4cm} C{2.0cm}@{}}
\toprule
$x$ & Working & $P(X=x)$ & Decimal \\
\midrule
0 & $1\cdot 5^{5}$   & $3125/7776$ & 0.4019 \\
1 & $5\cdot 5^{4}$   & $3125/7776$ & 0.4019 \\
2 & $10\cdot 5^{3}$  & $1250/7776$ & 0.1608 \\
3 & $10\cdot 5^{2}$  & $250/7776$  & 0.0322 \\
4 & $5\cdot 5^{1}$   & $25/7776$   & 0.0032 \\
5 & $1\cdot 5^{0}$   & $1/7776$    & 0.0001 \\
\midrule
\textbf{Total} & & \textbf{7776/7776} & \textbf{1.0000} \\
\bottomrule
\end{tabular}
\end{center}

\vcheck{$3125+3125+1250+250+25+1=7776=6^{5}$.}

\medskip
Note $E(X)=np=5\times\tfrac16=\tfrac56\approx 0.83$, consistent with the mass
piling up at $x=0$ and $x=1$.

\ans{$P(X=x)=\binom{5}{x}(1/6)^{x}(5/6)^{5-x}$, giving $0.4019$, $0.4019$,
$0.1608$, $0.0322$, $0.0032$, $0.0001$ for $x=0,\dots,5$.}
\end{solution}

%---------------------------------------------------------------------
\section{Three missing-value problems}

\begin{question}{Q.8, Q.9, Q.10}
\textbf{Q.8} The distribution of $X$ over $1,2,3,4,5$ is
$0.1,\;0.3,\;y,\;0.2,\;0.1$. Find $y$.

\smallskip
\textbf{Q.9} Find $k$ such that the values $0.01,\;0.25,\;0.40,\;k,\;0.04$ over
$X=2,3,4,5,6$ form a probability distribution.

\smallskip
\textbf{Q.10} For $X=0,2,4,6,8$ the probabilities are
$K,\;2K,\;4K,\;2K,\;K$. Rewrite the distribution after finding $K$.
\end{question}

\begin{solution}
All three are the same one-line idea: the probabilities must total 1.

\medskip
\textbf{Q.8.}
\[
0.1+0.3+y+0.2+0.1=1 \;\Longrightarrow\; 0.7+y=1 \;\Longrightarrow\; y=0.3 .
\]

\textbf{Q.9.}
\[
0.01+0.25+0.40+k+0.04=1 \;\Longrightarrow\; 0.70+k=1 \;\Longrightarrow\; k=0.30 .
\]

\textbf{Q.10.}
\[
K+2K+4K+2K+K=1 \;\Longrightarrow\; 10K=1 \;\Longrightarrow\; K=0.1 .
\]
The completed distribution is symmetric about $x=4$:

\begin{center}
\small\sffamily
\begin{tabular}{@{}l C{1.6cm} C{1.6cm} C{1.6cm} C{1.6cm} C{1.6cm}@{}}
\toprule
$x$      & 0   & 2   & 4   & 6   & 8   \\
$P(X=x)$ & 0.1 & 0.2 & 0.4 & 0.2 & 0.1 \\
\bottomrule
\end{tabular}
\end{center}

\vcheck{$0.1+0.2+0.4+0.2+0.1=1.0$; every value lies in $[0,1]$. Both conditions
hold. By symmetry $E(X)=4$.}

\ans{Q.8: $y=0.3$. \quad Q.9: $k=0.30$. \quad
Q.10: $K=0.1$, giving $0.1,\,0.2,\,0.4,\,0.2,\,0.1$.}
\end{solution}

\begin{pitfall}
After solving for the unknown, confirm it lies in $[0,1]$. A negative value
means the other probabilities were misread --- the answer is not ``$k=-0.2$'',
it is ``no such $k$ exists''.
\end{pitfall}
